\section{Aminoácidos}

\subsection{Aminoácidos canónicos}

Existe una gran cantidad de aminoácidos en la naturaleza, sin embargo, no todos son usados por esta para crear proteínas. Los que si son usados para tal tarea se conocen como aminoácidos canónicos, y se pueden observar en la \autoref{sec:aminoácidos_canónicos}.

\subsection{Efecto del pH en los aminoácidos}

Todo aminoácido tiene un grupo amino y un grupo carboxilo, y cada grupo puede existir en forma ácida o en forma básica dependiendo del pH de la solución en la que se encuentre en ese momento.

Las estructuras de las formas ácida y básica de cada grupo es:

    \begin{table}[H]
    \centering
    \caption{Grupos amino y carboxilo de los aminoácidos en sus formas ácidas y básicas}
    \label{fig:grupos_acidos-basicos}
    {\setstretch{1.25}
    \begin{tabular}{cc}
        \hline
        Formas básicas & Formas ácidas \\ \hline
        \begin{subfigure}{0.3\textwidth}
            \centering
            \vspace{0.25cm}
            \chemfig{H_3N^{+}-[:30]C-[:-30]C}
            \caption{Forma ácida del grupo amino}
            \label{fig:amino_acido}
        \end{subfigure}
        &
        \begin{subfigure}{0.3\textwidth}
            \centering
            \vspace{0.25cm}
            \chemfig{H_{2}N-[:30]C-[:-30]C}
            \caption{Forma básica del grupo amino}
            \label{fig:amino_básico}
        \end{subfigure}
        \\
        \begin{subfigure}{0.3\textwidth}
            \centering
            \vspace{0.25cm}
            \chemfig{HO-[:-30]C(=[:-90]O)-[:30]C}
            \caption{Forma ácida del grupo carboxilo}
            \label{fig:carboxilo_acido}
        \end{subfigure}
        &
        \begin{subfigure}{0.3\textwidth}
            \centering
            \vspace{0.25cm}
            \chemfig{O^{-}-[:-30]C(=[:-90]O)-[:30]C}
            \caption{Forma básica del grupo carboxilo}
            \label{fig:carboxilo_básico}
        \end{subfigure}
        \\ \hline
    \end{tabular}}
\end{table}

\definition{Anfoterismo}{Una sustancia anfótera es aquella que puede reaccionar ya sea como un ácido o como una base.}

Los aminoácidos disueltos en agua presentan un comportamiento anfótero, esto quiere decir que pueden ionizarse y comportarse tanto como ácido como base, esto dependiendo del pH del medio en que se encuentren.

\subsection{$pH$ y $pKa$}

El $pH$ puede afectar a la estructura de los aminoácidos. El valor de $pKa$ juega un papel importante en esto. Algunos valores se muestran en la \autoref{tab:pka_aminoácidos}

    \begin{table}
\centering
\caption{Lista de valores de pKa de algunos aminoácidos en base a su pH}
\label{tab:pka_aminoácidos}
\begin{threeparttable}
{\setstretch{2}
\begin{tabular}{cccc}\hline

    \textbf{Aminoácido} & \textbf{$pK_{a}$ de grupo carboxilo $COOH$} & \textbf{$pK_{a}$ de grupo amino $NH_{3}^{+}$} & $pK_{a}$ de cadena lateral\tnote{1} \\ \hline
    
    Ácido aspártico & 2.09 & 9.82 & 3.86 \\
    Ácido glutámico & 2.19 & 9.67 & 4.25 \\
    Alanina & 2.34 & 9.69 & - \\
    Arginina & 2.17 & 9.04 & 12.48 \\
    Asparagina & 2.02 & 8.84 & - \\
    Cisteína & 1.92 & 10.46 & 8.35 \\
    Fenilalanina & 2.16 & 9.18 & - \\
    Glicina & 2.34 & 9.60 & - \\
    Glutamina & 2.17 & 9.13 & - \\
    Histidina & 1.82 & 9.17 & 6.04 \\
    Isoleucina & 2.36 & 9.68 & - \\
    Leucina & 2.36 & 9.60 & - \\
    Lisina & 2.18 & 8.95 & 10.79 \\
    Metionina & 2.28 & 9.21 & - \\
    Prolina & 1.99 & 10.60 & - \\
    Serina & 2.21 & 9.15 & - \\
    Tirosina & 2.20 & 9.11 & 10.07 \\
    Treonina & 2.63 & 9.10 & - \\
    Trptófano & 2.38 & 9.39 & - \\
    Valina & 2.32 & 9.62 & - \\
    
    \hline
\end{tabular}
}
\begin{tablenotes} % ex: \tntote{1} ... \item[1]
    \item[1] Los valores que no aparecen (-) ...
\end{tablenotes}
\end{threeparttable}
\end{table}

\subsection{Utilidad del anfoterismo de los aminoácidos}

Debido a su comportamiento anfótero, los aminoácidos tienden a neutralizar las variaciones de pH del medio, ya que pueden comportarse como un ácido o una base, liberando o retirando protones del medio. Esta propiedad es importante en la homeostasis porque atenúan las variaciones de pH del medio, funcionando como un sistema tampón o amortiguador del pH.

Cada aminoácido tiene un pH en el que tiende a adoptar una forma dipolar neutra (o zwitterión), con tantas cargas positivas como negativas, que se denomina punto isoeléctrico.

\section{Péptidos}



\section{Proteínas}

En 1838, Berzelius escribió a Mulder y propuso la palabra "proteína" para estas nuevas moléculas, esto por ser la sustancia principal en la nutrición animal y en las cadenas tróficas.

\subsection{Síntesis de proteínas}

La síntesis proteica ocurre en los ribosomas, tomando de base estructural los 20 aminoácidos canónicos y traduciendo los ácidos nucleicos del material genético por medio del dogma central

\subsection{Funciones e importancia de las proteínas}

Entre las funciones de las proteínas podemos encontrar:

\begin{itemize}
    \item Estructural
    \item Enzimática
    \item Comunicación celular
    \item Transporte
    \item Movilidad
\end{itemize}

Las proteínas son de importancia por varias razones:

\begin{itemize}
    \item Son abundantes y diversas
    \item Conforman 50\% del peso seco celular
    \item Tienen gran importancia estructural
    \item Son muy versátiles
\end{itemize}